% GENERATED FILE. Edit canonical lessons, then run npm run generate:books.
% canonical-source-hash: fnv1a64:fb350029c6dce76b
% canonical-lessons: IT-C185-muscolo, IT-C185-pollice, IT-C185-ferita, IT-C185-cicatrice, IT-C185-raffreddore

\chapter{Muscle, Thumb, Wound, Scar, Cold}
\label{ch:it-muscle-thumb-wound-scar-cold}
\hlchaptermodality{185}

\begin{chapteropening}
\textbf{By the end of this chapter:} \emph{I can say a muscle, a thumb, a wound, a scar and a cold.}

Complete the last lesson of chapter 185: I can say a muscle, a thumb, a wound, a scar and a cold.
\end{chapteropening}

\section[il muscolo]{il muscolo — a muscle}

\label{lesson:IT-C185-muscolo}



\begin{quote}
Before the new one: say the Italian for to suspect, then the Italian for to squeeze.
\end{quote}

\subsection*{You'll want to know: il muscolo}
\textbf{il muscolo} — \textquotedblleft{}a muscle\textquotedblright{}.

\begin{cousinweb}
\textbf{muscolo} $\leftarrow$ Latin \emph{musculus}, \textquotedblleft{}a little mouse; a muscle\textquotedblright{} — a learned borrowing: a flexing muscle looked like a mouse under the skin. (EtymDB 2.1 records the Latin word.)
\end{cousinweb}

\begin{grammarlens}[title={What you've built}]
One more word for this chapter.
\end{grammarlens}

\subsection*{Your turn}
\begin{itemize}
\raggedright
  \item \emph{Say it:} \emph{il muscolo}
  \item \emph{Say it:} \emph{il muscolo}, once more
  \item \emph{Say it:} say \emph{il muscolo}
  \item \emph{Recall:} say the Italian for to suspect, then the Italian for to squeeze, then say \emph{il muscolo} again
\end{itemize}

\begin{tcolorbox}[breakable,colback=teal!4,colframe=teal!35!black,title={Before you move on}]
What does il muscolo mean? (\textquotedblleft{}A muscle\textquotedblright{}.) Say it once more.
\end{tcolorbox}
\section[il pollice]{il pollice — a thumb}

\label{lesson:IT-C185-pollice}



\begin{quote}
Before the new one: say the Italian for to squeeze, then the Italian for a muscle.
\end{quote}

\subsection*{You'll want to know: il pollice}
\textbf{il pollice} — \textquotedblleft{}a thumb\textquotedblright{}.

\begin{cousinweb}
\textbf{pollice} $\leftarrow$ Latin \emph{pollex}, \textquotedblleft{}a thumb\textquotedblright{} (\emph{pollicem}) — inherited. (EtymDB 2.1)
\end{cousinweb}

\begin{grammarlens}[title={What you've built}]
One more word for this chapter.
\end{grammarlens}

\subsection*{Your turn}
\begin{itemize}
\raggedright
  \item \emph{Say it:} \emph{il pollice}
  \item \emph{Say it:} \emph{il pollice}, once more
  \item \emph{Say it:} say \emph{il pollice}
  \item \emph{Recall:} say the Italian for to squeeze, then the Italian for a muscle, then say \emph{il pollice} again
\end{itemize}

\begin{tcolorbox}[breakable,colback=teal!4,colframe=teal!35!black,title={Before you move on}]
What does il pollice mean? (\textquotedblleft{}A thumb\textquotedblright{}.) Say it once more.
\end{tcolorbox}
\section[la ferita]{la ferita — a wound}

\label{lesson:IT-C185-ferita}



\begin{quote}
Before the new one: say the Italian for a muscle, then the Italian for a thumb.
\end{quote}

\subsection*{You'll want to know: la ferita}
\textbf{la ferita} — \textquotedblleft{}a wound\textquotedblright{}.

\begin{cousinweb}
\textbf{ferita} is the past participle of \emph{ferire}, \textquotedblleft{}to wound\textquotedblright{}, used as a noun: a thing struck. \emph{ferire} $\leftarrow$ Latin \emph{ferire}, \textquotedblleft{}to strike\textquotedblright{}. (EtymDB 2.1)
\end{cousinweb}

\begin{grammarlens}[title={What you've built}]
One more word for this chapter.
\end{grammarlens}

\subsection*{Your turn}
\begin{itemize}
\raggedright
  \item \emph{Say it:} \emph{la ferita}
  \item \emph{Say it:} \emph{la ferita}, once more
  \item \emph{Say it:} say \emph{la ferita}
  \item \emph{Recall:} say the Italian for a muscle, then the Italian for a thumb, then say \emph{la ferita} again
\end{itemize}

\begin{tcolorbox}[breakable,colback=teal!4,colframe=teal!35!black,title={Before you move on}]
What does la ferita mean? (\textquotedblleft{}A wound\textquotedblright{}.) Say it once more.
\end{tcolorbox}
\section[la cicatrice]{la cicatrice — a scar}

\label{lesson:IT-C185-cicatrice}



\begin{quote}
Before the new one: say the Italian for a thumb, then the Italian for a wound.
\end{quote}

\subsection*{You'll want to know: la cicatrice}
\textbf{la cicatrice} — \textquotedblleft{}a scar\textquotedblright{}.

\begin{cousinweb}
\textbf{cicatrice} $\leftarrow$ Latin \emph{cicatrix}, \textquotedblleft{}a scar\textquotedblright{} — inherited. (EtymDB 2.1)
\end{cousinweb}

\begin{grammarlens}[title={What you've built}]
One more word for this chapter.
\end{grammarlens}

\subsection*{Your turn}
\begin{itemize}
\raggedright
  \item \emph{Say it:} \emph{la cicatrice}
  \item \emph{Say it:} \emph{la cicatrice}, once more
  \item \emph{Say it:} say \emph{la cicatrice}
  \item \emph{Recall:} say the Italian for a thumb, then the Italian for a wound, then say \emph{la cicatrice} again
\end{itemize}

\begin{tcolorbox}[breakable,colback=teal!4,colframe=teal!35!black,title={Before you move on}]
What does la cicatrice mean? (\textquotedblleft{}A scar\textquotedblright{}.) Say it once more.
\end{tcolorbox}
\section[il raffreddore]{il raffreddore — a cold}

\label{lesson:IT-C185-raffreddore}



\begin{quote}
Before the new one: say the Italian for a wound, then the Italian for a scar.
\end{quote}

\subsection*{You'll want to know: il raffreddore}
\textbf{il raffreddore} — \textquotedblleft{}a cold\textquotedblright{}.

\begin{cousinweb}
\textbf{raffreddore} — built on \emph{freddo}, \textquotedblleft{}cold\textquotedblright{}, with \emph{ra-} and \emph{-ore}: a chill. \emph{freddo} goes back to Latin \emph{frigidus}, \textquotedblleft{}cold\textquotedblright{}. EtymDB 2.1 has no row for \emph{raffreddore}.
\end{cousinweb}

\begin{grammarlens}[title={What you've built}]
One more word for this chapter.
\end{grammarlens}

\subsection*{Your turn}
\begin{itemize}
\raggedright
  \item \emph{Say it:} \emph{il raffreddore}
  \item \emph{Say it:} \emph{il raffreddore}, once more
  \item \emph{Say it:} say \emph{il raffreddore}
  \item \emph{Recall:} say the Italian for a wound, then the Italian for a scar, then say \emph{il raffreddore} again
\end{itemize}

\begin{tcolorbox}[breakable,colback=teal!4,colframe=teal!35!black,title={Before you move on}]
What does il raffreddore mean? (\textquotedblleft{}A cold\textquotedblright{}.) Say it once more.
\end{tcolorbox}
